Based on the observations of the previous sections, we now define a standard for alignments.
Because many of the alignment types described in \autoref{sec:crosslib:alignments:types} are very difficult to handle rigorously and additional alignment types may be discovered in the future, we opt for a very simple and flexible definition.

Concretely, we use the following formal grammar for collections of alignments:
\begin{center}
\begin{tabular}{l@{\tb}c@{\tb}l}
	Collection  & \bbc  & $\rep{\bnfbracket{\mathrm{Comment} \bnfalt \mathrm{NSDef} \bnfalt \mathrm{Alignment}}}$\\
	Comment & \bbc & // String \\
	NSDef & \bbc  &\texttt{namespace} String URI\\
	Alignment & \bbc & URI URI $\rep{\bnfbracket{\mathrm{String} = "\mathrm{String}"}}$
\end{tabular}
\end{center}
Our definition aims at practicality, especially considering the typical case where researchers exchange and manipulate large collections of alignments.
Therefore, our grammar allows for comments and for the introduction of short namespace definitions that abbreviate long namespaces.
Our grammar represents each individual alignment as a pair of two URIs with arbitrary additional data stored as a list of key-value pairs.

The additional data in alignments makes our standard extensible: any user can standardize individual keys in order to define specific types of alignments.
For example, for alignments up to argument order, we can add a key for giving the argument order.
Moreover, this can be used to annotate metadata such as provenance or system versions.

Below we standardize some individual keys and use them to implement the most important alignment types from \autoref{sec:crosslib:alignments:types}.
In all definitions below, we assume that $a_1$ and $a_2$ are the aligned symbols.

\begin{definition}
The key \cn{direction} has the possible values \cn{forward}, \cn{backward}, or \cn{both}.
Its presence legitimizes, respectively, the translation that replaces every occurrence of $a_1$ with $a_2$, its inverse, or both.
\end{definition}
Alignments with \cn{direction} key subsume the alignment types of perfect alignments (where the direction is \cn{both}) and the unidirectional types of alignment up to totality of functions or up to associativity, and alignment for certain arguments.
The absence of this key indicates those alignment types where no symbol-to-symbol translation is possible, in particular opaque alignments.

\begin{definition}\label{def:crosslib:alignments:arguments}
The key \cn{arguments} has values of the form $(r_1,s_1)\ldots(r_k,s_k)$ where the $r_i$ and $s_i$ are natural numbers.
Its presence legitimizes the translation of $a_1(x_1,\ldots,x_m)$ to $a_2(y_1,\ldots,y_n)$ where each $y_k$ is defined by
\begin{compactitem}
 \item if $k=s_i$ for some $i$: the recursive translation of $x_{r_i}$ 
 \item otherwise: inferred from the context.
\end{compactitem}
\end{definition}
Alignments with \cn{arguments} key subsume the alignment types of alignments up to argument order and of alignment up to determined arguments.

\begin{example}
We obtain the following argument alignments for some of the examples from \autoref{sec:crosslib:alignments:types}:
\[\begin{array}{c}
    \cn{Nat}_1 \; \cn{Nat}_2 \; \cn{direction}="\cn{both}"\\
    \cn{eq}_1 \; \cn{eq}_2 \; \cn{arguments}="(1,2)(2,3)"\\
    \cn{contains}_1 \; \cn{in}_2 \; \cn{arguments}="(1,1)(2,3)(3,2)"
  \end{array}
  \]
\end{example}

%Another possible case is when two symbols are aligned in the same library.
%Although can be recovered by comparing their URIs, it is made explicit by the following tag.
%\begin{definition}
%  An \textbf{intracorporal alignment} uses the key \cn{self} with the unique value \cn{true}.
%\end{definition}

Finally, we standardize a key for probabilistic alignments:
\begin{definition}
The key \cn{similarity} has values that are real numbers in $[0;1]$.
If used together with other keys like \cn{direction} and \cn{arguments}, it represents a certainty score for the correctness of the corresponding translation.
If absent, its value is $1$ indicating perfect certainty.
\end{definition}
%When working with a large set of alignments between a fixed pair of corpora, this key describes the similarity score is most indicative when comparing alignments found in the same pair of corpora.
%This is the case for many of the alignments we have found so far.

%%% Local Variables:
%%% mode: latex
%%% TeX-master: "paper-cicm2"
%%% End:

%  LocalWords:  bnfbracket mathrm bnfalt mathrm bnfalt mathrm standardize cn cn cn cn
%  LocalWords:  ldots ldots
